\documentclass[10pt]{amsart}

\pdfinfoomitdate=1
\pdftrailerid{}
\pdfsuppressptexinfo=15

\usepackage[T1]{fontenc}
\usepackage{lmodern}
\usepackage[a4paper,margin=25mm]{geometry}
\usepackage{amsmath,amssymb,mathtools}
\usepackage{enumitem}
\usepackage{microtype}
\usepackage[hidelinks]{hyperref}

\setlist{nosep,leftmargin=*}
\setlength{\parskip}{2pt}
\setlength{\parindent}{15pt}
\emergencystretch=2em

\newtheorem{theorem}{Theorem}[section]
\newtheorem{proposition}[theorem]{Proposition}
\newtheorem{lemma}[theorem]{Lemma}
\newtheorem{corollary}[theorem]{Corollary}
\theoremstyle{definition}
\newtheorem{definition}[theorem]{Definition}
\newtheorem{remark}[theorem]{Remark}

\newcommand{\Par}{\mathcal P}
\newcommand{\Img}{\mathcal I}
\newcommand{\Fix}{\operatorname{Fix}}
\newcommand{\ellp}{\ell}

\title[Principal-hook partition dynamics]{Principal-hook regrouping dynamics on integer partitions}
\author{Anonymous}
\date{}

\begin{document}

\begin{abstract}
For a partition of a fixed positive integer, regroup the Ferrers diagram by
its principal diagonal hooks and iterate the resulting partition.  The
principal-hook partition and first-hook identity, the adjacent-gap image and
product fibre formula, and the one-step diagonal-hook formalism are all
treated as owned background and receive zero credit.  Our single main theorem
is an exact increment identity for the first adjacent gap, from which we
deduce the sharp maximum transient depth $\lfloor n/2\rfloor$.  As explicit
low-credit corollaries, the one-row partition is globally absorbing, the
owned fibre weights give a depth-state-weighted layer transport identity,
conjugation has one entrance-time exception, and each fixed-weight system has
zeta function $(1-z)^{-1}$.  All statements include $n=1,2$.
\end{abstract}

\maketitle

\section{Introduction and scope}

Let $\Par(n)$ be the finite set of integer partitions of $n\geq 1$.  We study
the deterministic self-map obtained by replacing a Ferrers diagram with the
list of its principal diagonal-hook lengths.  Principal hooks and Frobenius
coordinates are standard tools in partition theory \cite{Andrews1998}.  We
make the ownership subtraction item by item.  Gutschwager explicitly defines
the principal-hook length partition and records its first part
$\lambda_1+\ell(\lambda)-1$ \cite{Gutschwager2011}; Goupil records the
adjacent-gap image and exact product fibre weight \cite{Goupil2009}; and
Chern--Yee give direct prior work on diagonal-hook data and an involution
preserving every diagonal-hook length \cite{ChernYee2022}.  The map object,
first-hook identity, one-step image, fibre product, and standard
diagonal-hook symmetries therefore receive zero credit here.
Proposition~\ref{prop:classical} is reproduced only to keep the later dynamics
self-contained.

After this subtraction, the present proof package has one main theorem: a
direct Ferrers-diagram calculation gives the exact increment of the first
adjacent gap, a pointwise depth bound, and the sharp global depth.  Global
absorption from the owned first-hook identity, fibre-weighted layer transport,
the conjugation timing exception, the periodic census, and zeta are labelled
low-credit corollaries rather than co-equal headline results.  A bounded owner
search located no exact temporal owner, but this negative search is not
novelty or priority evidence.

\paragraph{Firewall.}
The state space here consists of unlabelled integer partitions of a fixed
weight, and the update regroups cells into diagonal hooks.  It is not the
labelled set-partition lattice, and it uses neither a cyclic action nor a
lattice join.  Consequently none of the cyclic-shift join engine, Bell-number
basin census, or M\"obius-lattice mechanism of the separate P110 system is
invoked or claimed.  External novelty, priority, and dissemination status for
the results below remain \textsc{hold}.

\section{Principal hooks and owned one-step inputs}

Write a partition as
$\lambda=(\lambda _1,\ldots,\lambda _{\ellp})$ with
$\lambda_1\geq\cdots\geq\lambda_{\ellp}>0$, and write $\lambda'$ for its
conjugate.  Its Durfee size is
\[
 d=d(\lambda)=\max\{i:\lambda_i\geq i\}.
\]
The Frobenius arms and legs are
\[
 \alpha_i=\lambda_i-i,
 \qquad
 \beta_i=\lambda'_i-i
 \qquad(1\leq i\leq d).
\]
Both sequences are strictly decreasing and nonnegative.  The $i$th principal
diagonal hook has length
\[
 h_i=\alpha_i+\beta_i+1
     =\lambda_i+\lambda'_i-2i+1.
\]
These hooks partition the Ferrers diagram: a cell $(r,c)$ belongs to the hook
rooted at $(\min\{r,c\},\min\{r,c\})$.  Hence $\sum_i h_i=|\lambda|$, while
$h_i-h_{i+1}\geq2$.  Thus
\[
 H(\lambda)=(h_1,\ldots,h_d)
\]
defines a self-map of $\Par(n)$.  This principal-hook partition is a standard
object, explicitly denoted $hl(\lambda)$ by Gutschwager
\cite{Gutschwager2011}; its use as a one-step map receives zero credit.

Let
\[
 \Img(n)=\{(h_1,\ldots,h_r)\in\Par(n):
               h_i-h_{i+1}\geq2\text{ for }1\leq i<r\}.
\]
For a one-part partition the gap condition is vacuous.

\begin{proposition}[Previously recorded one-step result; zero credit]
\label{prop:classical}
The image of $H:\Par(n)\to\Par(n)$ is exactly $\Img(n)$.  Moreover, if
$h=(h_1,\ldots,h_r)\in\Img(n)$, then
\begin{equation}
 \label{eq:fibre}
 w(h):=\#H^{-1}(h)
 =h_r\prod_{i=1}^{r-1}(h_i-h_{i+1}-1).
\end{equation}
The empty product in \eqref{eq:fibre} is $1$.
\end{proposition}

\begin{proof}
The preceding calculation puts every image in $\Img(n)$.  Conversely,
Frobenius coordinates over a prescribed $h$ are precisely pairs of strictly
decreasing nonnegative sequences $(\alpha_i)$ and $(\beta_i)$ satisfying
$\alpha_i+\beta_i=h_i-1$.  There are $h_r$ choices for
$\alpha_r\in\{0,\ldots,h_r-1\}$, after which $\beta_r$ is fixed.  For
$i<r$, put
\[
 x_i=\alpha_i-\alpha_{i+1}-1,
 \qquad
 y_i=\beta_i-\beta_{i+1}-1.
\]
Then $x_i,y_i\geq0$ and
$x_i+y_i=h_i-h_{i+1}-2$, giving
$h_i-h_{i+1}-1$ choices.  These choices are independent and reconstruct the
two Frobenius sequences, hence a unique partition.  This proves both
surjectivity onto $\Img(n)$ and \eqref{eq:fibre}.  The argument and product
were previously recorded by Goupil \cite{Goupil2009}.
\end{proof}

\section{The gap clock and sharp depth}

A fixed point $q$ of a map on a finite set will be called
\emph{globally absorbing}, or \emph{globally attracting in finite time}, if
every state reaches $q$ after finitely many iterates.  This is the only
meaning of global attraction used below.

For $\lambda\in\Par(n)$ define its entrance time
\[
 \tau(\lambda)=\min\{t\geq0:H^t(\lambda)=(n)\}.
\]
The next result first shows that this number is always defined.

\begin{corollary}[Owned first-hook identity and global absorption; low credit]
\label{cor:absorption}
For every $\lambda\in\Par(n)$,
\begin{equation}
 \label{eq:firstpart}
 (H\lambda)_1=\lambda_1+\ellp(\lambda)-1.
\end{equation}
Unless $\lambda=(n)$, the right-hand side is strictly larger than
$\lambda_1$.  Consequently $(n)$ is a globally absorbing fixed point and is
the unique fixed point and the unique periodic point of $H$ on $\Par(n)$.
The identity \eqref{eq:firstpart} itself is owned and receives zero credit.
\end{corollary}

\begin{proof}
The first diagonal hook consists of the first row and first column with their
common corner counted once, which gives the identity already recorded in
\cite{Gutschwager2011}.  A non-one-row partition has
$\ellp(\lambda)\geq2$, so its first part strictly increases.  Since
$\Par(n)$ is finite and no first part exceeds $n$, every orbit reaches
$(n)$.  The same strict increase excludes any other fixed or periodic point.
\end{proof}

Pad a one-part partition by $\lambda_2=0$ and set
\[
 g(\lambda)=\lambda_1-\lambda_2,
 \qquad
 m_1(\lambda)=\#\{i:\lambda_i=1\}.
\]

\begin{theorem}[Main theorem: exact increment and sharp depth]
\label{thm:gapdepth}
Let $\lambda\in\Par(n)$ be nonterminal.
\begin{enumerate}[label=\textup{(\roman*)}]
\item If $d(\lambda)\geq2$, then
\begin{equation}
 \label{eq:gapinc}
 g(H\lambda)-g(\lambda)
 =\ellp(\lambda)-\lambda'_2+2
 =2+m_1(\lambda)\geq2.
\end{equation}
\item If $d(\lambda)=1$, then
$\lambda=(a,1^b)$ for some $b\geq1$, $H\lambda=(n)$, and
\begin{equation}
 \label{eq:durfeeone}
 g(H\lambda)-g(\lambda)=b+1\geq2.
\end{equation}
\end{enumerate}
For every $\lambda\in\Par(n)$,
\begin{equation}
 \label{eq:pointbound}
 \tau(\lambda)\leq
 \left\lfloor\frac{n-g(\lambda)}2\right\rfloor
 \leq\left\lfloor\frac n2\right\rfloor.
\end{equation}
The global maximum is sharp:
\begin{equation}
 \label{eq:maxdepth}
 \max_{\lambda\vdash n}\tau(\lambda)=\left\lfloor\frac n2\right\rfloor.
\end{equation}
It is attained by
$\lambda^{\star}=(\lceil n/2\rceil,\lfloor n/2\rfloor)$, with the zero
second part omitted when $n=1$.
\end{theorem}

\begin{proof}
Suppose first that $d\geq2$.  The first two parts of $H\lambda$ are
\[
 h_1=\lambda_1+\ellp(\lambda)-1,
 \qquad
 h_2=\lambda_2+\lambda'_2-3.
\]
Subtracting and then subtracting $g(\lambda)$ proves the first equality in
\eqref{eq:gapinc}.  Because $\lambda'_2$ counts the parts of $\lambda$ that
are at least two, $\ellp(\lambda)-\lambda'_2=m_1(\lambda)$, proving the
second equality.

If $d=1$, every part below the first is $1$.  Nonterminality gives
$\lambda=(a,1^b)$ with $b\geq1$.  Its sole principal hook is the entire
diagram, so $H\lambda=(n)$; since $n=a+b$ and $g(\lambda)=a-1$,
\eqref{eq:durfeeone} follows.

At each of the $\tau(\lambda)$ nonterminal steps the gap therefore rises by
at least two, whereas the terminal gap is $g((n))=n$.  This gives
$2\tau(\lambda)\leq n-g(\lambda)$ and hence \eqref{eq:pointbound}; it also
holds trivially at the terminal state.

It remains to attain the bound.  For $n=1$, the state $(1)$ is terminal and
has depth zero.  If $a\geq b\geq2$, a direct two-row Ferrers
calculation gives
\begin{equation}
 \label{eq:tworow}
 H(a,b)=(a+1,b-1),
\end{equation}
whereas $H(a,1)=(a+1)$.  For $n\geq2$, start from
$a=\lceil n/2\rceil$ and $b=\lfloor n/2\rfloor$, so $b\geq1$.  Equation
\eqref{eq:tworow} applies $b-1$ times (zero times when $b=1$), ending at
$(n-1,1)$; the final hook step $H(n-1,1)=(n)$ gives $b$ steps in total.
Thus the upper bound is attained and \eqref{eq:maxdepth} follows.
\end{proof}

\begin{remark}[The two smallest weights]
For $n=1$, $\Par(1)=\{(1)\}$ and the maximum depth is $0$.  For $n=2$,
$H(1,1)=(2)$, so the two depths are $0$ and $1$.  These agree with
Theorem~\ref{thm:gapdepth}; no exceptional empty state is introduced.
\end{remark}

\section{Low-credit consequences}

For $t\geq0$ let
\[
 A_t(n)=\#\{\lambda\in\Par(n):\tau(\lambda)=t\}.
\]

\begin{corollary}[Fibre-weighted layer transport; low credit]
\label{cor:layers}
For every $n\geq1$,
\begin{equation}
 \label{eq:initiallayers}
 A_0(n)=1,
 \qquad
 A_1(n)=n-1.
\end{equation}
For every $t\geq2$,
\begin{equation}
 \label{eq:layerrecurrence}
 A_t(n)=
 \sum_{\substack{h=(h_1,\ldots,h_r)\in\Img(n)\\
                   \tau(h)=t-1}}
 h_r\prod_{i=1}^{r-1}(h_i-h_{i+1}-1).
\end{equation}
An empty sum is $0$.  In particular, $A_t(n)=0$ for
$t>\lfloor n/2\rfloor$ and $\sum_tA_t(n)=|\Par(n)|$.
This is a depth-state-weighted transport identity over $\Img(n)$, not a
closed scalar recurrence in the numbers $A_t(n)$ alone.
\end{corollary}

\begin{proof}
Corollary~\ref{cor:absorption} makes $(n)$ the only depth-zero state.  The fibre
over $(n)$ has size $w((n))=n$ by Proposition~\ref{prop:classical}; exactly
one member of that fibre is $(n)$ itself.  This proves
\eqref{eq:initiallayers}, including $A_1(1)=0$.

If $t\geq2$, a state $\lambda$ has depth $t$ exactly when its first image
$h=H\lambda$ has depth $t-1$.  Every possible $h$ belongs to $\Img(n)$, and
the fibres over distinct $h$ are disjoint.  Summing the owned fibre
weight \eqref{eq:fibre} gives \eqref{eq:layerrecurrence}.  The vanishing range
follows from Theorem~\ref{thm:gapdepth}, and the mass identity follows because
the depth layers partition the finite state space.
\end{proof}

\begin{corollary}[Conjugation symmetry and timing consequence; low credit]
\label{cor:conjugation}
For every $\lambda\in\Par(n)$,
\begin{equation}
 \label{eq:conjugation}
 H(\lambda)=H(\lambda').
\end{equation}
Consequently $H^t(\lambda)=H^t(\lambda')$ for every $t\geq1$.  Entrance
depth is also conjugation-invariant except for one explicit pair: if $n>1$,
then
\[
 \tau((n))=0,
 \qquad
 \tau((1^n))=1.
\]
There are no other failures of $\tau(\lambda)=\tau(\lambda')$.
The one-step diagonal-hook symmetry in \eqref{eq:conjugation} is standard and
receives zero credit; only its stated timing consequence is used here.
\end{corollary}

\begin{proof}
Conjugation swaps the Frobenius arm and leg sequences, leaving every sum
$\alpha_i+\beta_i+1$ unchanged.  This proves \eqref{eq:conjugation} and the
iterate statement.  If neither member of a conjugate pair is terminal, then
both depths equal $1+\tau(H\lambda)$.  The only time exactly one member is
terminal is the pair $(n),(1^n)$ for $n>1$; its displayed depths follow from
the single-hook calculation.  At $n=1$ the two partitions coincide.
\end{proof}

\begin{corollary}[Fixed-weight periodic census and zeta; low credit]
\label{cor:zeta}
For each fixed $n\geq1$, write $H_n:\Par(n)\to\Par(n)$ for the map above.
For every $m\geq1$,
\[
 \#\Fix(H_n^m)=1.
\]
Hence the Artin--Mazur zeta function of this fixed finite system is
\begin{equation}
 \label{eq:zeta}
 \zeta_{H_n}(z)
 =\exp\!\left(\sum_{m\geq1}\frac{\#\Fix(H_n^m)}m z^m\right)
 =\frac1{1-z}
\end{equation}
as a formal power series.  The weight $n$ is fixed before forming the zeta
function; no zeta function on the disjoint union of all weights is asserted.
\end{corollary}

\begin{proof}
Every point fixed by an iterate is periodic.  Corollary~\ref{cor:absorption}
excludes every periodic point except $(n)$, which is fixed by all iterates.
Thus every fixed-point count is one.  Substitution into the definition gives
$\exp(\sum_{m\geq1}z^m/m)=\exp(-\log(1-z))$, proving
\eqref{eq:zeta}.
\end{proof}

\section{Two proof routes and falsification controls}

The argument deliberately uses two distinct mechanisms.  The
\emph{Frobenius/fibre route} decomposes each hook length into a strict arm and
leg pair; it verifies the previously recorded one-step characterization and
product weight, which remain zero-credit background.  The
\emph{Ferrers/gap route} reads the first two diagonal hooks directly from row
and column lengths; its substantive output is the main exact increment
\eqref{eq:gapinc} and sharp depth theorem.  Global absorption and the periodic
census are elementary low-credit consequences of the owned first-hook
identity.  The layer transport joins the two routes but remains a low-credit,
nonclosed identity.

Several tempting strengthenings are false.  The map is not a projection:
$H(2,2)=(3,1)$ and $H^2(2,2)=(4)$.  Corollary~\ref{cor:conjugation} cannot
be strengthened to unconditional depth invariance, already at
$(2)$ and $(1,1)$.  Nor does a simple rectangular boundary describe all
deepest states: for example $(4,4,4,4)\vdash16$ has depth $7$, below the
maximum $8$.  These examples, the formulas above, and all partitions through
$n=40$ are checked by the exact standard-library verifier distributed with
the manuscript.  Computation is used only as a falsification and regression
control; all mathematical statements have independent proofs above.

\bibliographystyle{amsplain}
\bibliography{references}

\end{document}
